% !TEX program = xelatex
% =============================================================================
%  《量子场论与 Jones 多项式》读书笔记 —— 导言区
%  编译：xelatex（两遍）；中文字体：fandol
% =============================================================================
\documentclass[11pt,a4paper]{article}

\usepackage[UTF8, fontset=fandol, heading=true]{ctex}
\usepackage{amsmath}
\usepackage{amssymb}
\usepackage{mathtools}
\usepackage{bm}
\usepackage[a4paper,left=2.6cm,right=2.6cm,top=2.7cm,bottom=2.7cm]{geometry}
\usepackage{graphicx}
\usepackage{xcolor}
\usepackage{tcolorbox}
\tcbuselibrary{breakable, skins}
\usepackage{enumitem}
\usepackage{booktabs}
\usepackage{longtable}
\usepackage{array}
\usepackage[font=small,labelfont=bf,labelsep=quad]{caption}
\usepackage[colorlinks=true,
            linkcolor=blue!45!black,
            citecolor=blue!45!black,
            urlcolor=blue!45!black]{hyperref}

\linespread{1.15}
\setlength{\emergencystretch}{3em}
\allowdisplaybreaks[3]

\renewcommand{\figurename}{图}
\renewcommand{\tablename}{表}
\renewcommand{\refname}{参考文献}
\renewcommand{\contentsname}{目录}

\hypersetup{
  pdftitle={量子场论与 Jones 多项式：读书笔记},
  pdfauthor={WIT-1989-JONES 研讨班},
  pdfsubject={Chern-Simons theory, Jones polynomial, geometric quantization},
  pdfkeywords={Chern-Simons, Jones polynomial, conformal blocks, Verlinde formula},
  bookmarksnumbered=true
}

\newcommand{\ii}{\mathrm{i}}
\DeclareMathOperator{\Tr}{Tr}
\DeclareMathOperator{\Hom}{Hom}
\DeclareMathOperator{\Inv}{Inv}
\newcommand{\dd}{\mathrm{d}}
\newcommand{\CC}{\mathbb{C}}
\newcommand{\ZZ}{\mathbb{Z}}
\newcommand{\cM}{\mathcal{M}}
\newcommand{\cL}{\mathcal{L}}
\newcommand{\cH}{\mathcal{H}}
\newcommand{\cO}{\mathcal{O}}
\newcommand{\cD}{\mathcal{D}}
\newcommand{\cN}{\mathcal{N}}

% ---- 三种断言标记：原文 / 我们推导 / 我们猜测 -------------------------------
\newcommand{\src}[1]{\textnormal{\footnotesize[原文 p.#1]}}
\newcommand{\ours}{\textnormal{\footnotesize[我们推导]}}
\newcommand{\guess}{\textnormal{\footnotesize[我们猜测]}}

% ---- 章末「边界」框 ---------------------------------------------------------
\newenvironment{boundary}{%
    \par\medskip
    \begin{tcolorbox}[enhanced, breakable,
        colback=orange!4!white, colframe=orange!55!black,
        boxrule=0.5pt, arc=2pt, left=6pt, right=6pt, top=4pt, bottom=4pt,
        title={\small\textbf{本章边界}\quad{\footnotesize（哪些已证、哪些未证）}},
        fonttitle=\bfseries, coltitle=black, colbacktitle=orange!12!white]
}{%
    \end{tcolorbox}
    \par\medskip
}

% ---- 关键结论框 -------------------------------------------------------------
\newenvironment{keypoint}{%
    \par\medskip
    \begin{tcolorbox}[enhanced, breakable,
        colback=blue!3!white, colframe=blue!40!black,
        boxrule=0.5pt, arc=2pt, left=6pt, right=6pt, top=4pt, bottom=4pt]
}{%
    \end{tcolorbox}
    \par\medskip
}

% ---- 排版容差 ----
\tolerance=3000
\emergencystretch=3em
\hbadness=10000


\title{\bfseries 量子场论与 Jones 多项式\\[4pt]
\large 一篇 1989 年论文的逐章导读}
\author{WIT-1989-JONES 研讨班\\[2pt]
\small 原著：Edward Witten, \emph{Quantum Field Theory and the Jones Polynomial},\\
\small \emph{Commun. Math. Phys.} \textbf{121} (1989) 351--399}
\date{2026 年}

\begin{document}

\maketitle

\begin{abstract}
\noindent
这篇笔记要回答一个问题：Witten 在 1989 年做了什么，使 Jones 多项式第一次有了内蕴的三维定义？答案的主线只有一句：把一个不含度规的三维作用量（Chern--Simons 作用量）量子化，沿曲面切开时，量子 Hilbert 空间恰好是二维共形场论的共形块空间；纽结作为 Wilson 线进入这个理论，于是纽结不变量变成三维流形不变量。笔记按九个问题展开，每章给出论证链、关键公式的起点与终点，并在章末标出"作者证明了什么、我们补了什么、哪些仍是猜想"。
\end{abstract}

\tableofcontents
\newpage

\setcounter{section}{-1}
\section{五分钟版：这篇文章到底做了什么}

如果你只有五分钟，读这一节。

\begin{keypoint}
\textbf{一句话}：Witten 把三维 Chern--Simons 规范理论量子化，得到的量子态空间随曲面变化构成一个模函子；二维共形场论的共形块正是这个模函子的取值；纽结作为 Wilson 线插入后，配分函数就是纽结多项式，而且它对三维流形的改变（手术）给出不变量。
\end{keypoint}

论证分四步，每步的“钥匙”如下：

\begin{enumerate}[leftmargin=2em]
\item \textbf{写出一个不含度规的作用量。}
Chern--Simons 作用量 $L=\frac{k}{4\pi}\int_M\Tr(A\wedge \dd A+\tfrac23 A\wedge A\wedge A)$ 只依赖定向，不依赖度规 \src{353--354}。度规不在作用量里，所以量子化得到的量自动是拓扑不变量。这是“广义协变”的一种实现，与广义相对论“对度规积分”不同。

\item \textbf{沿曲面切开，读出相空间。}
把三维流形写成 $\Sigma\times\mathbb{R}$，作用量对时间是一阶的，给出泊松括号 $\{A_i^a(x),A_j^b(y)\}=\frac{4\pi}{k}\epsilon_{ij}\delta^{ab}\delta^2(x-y)$ \src{367}。高斯定律约束 $F=0$ 是第一类约束，辛商之后相空间是“$\Sigma$ 上平坦联络模去规范变换”的有限维空间 $\cM$ \src{367}。有限维，所以量子 Hilbert 空间有限维。

\item \textbf{几何量子化：从相空间到全纯截面。}
$\cM$ 在选一个复结构 $J$ 后是 K\"ahler 簇；预量子化线丛取行列式线丛 $\cL$，量子态空间是 $\cH_\Sigma^{(J)}=H^0(\cM_J,\cL^{\otimes k})$ \src{368}。$J$ 变时这些空间构成模空间上的射影平坦向量丛，“平坦”来自共形反常的单值性，单值性正是映射类群的作用 \src{369--371}。

\item \textbf{纽结进来，共形块对上。}
Wilson 线 $W_R(C)=\Tr_R P\exp\int_C A$ 不需要度规 \src{355}。穿刺点上的高斯定律是算符值的约束，Witten 的解法不是解它，而是把表示 $R$ 也换成量子对象：用 Borel--Weil--Bott 定理把 $R$ 实现为旗流形 $G/T$ 上辛形式 $\omega_R$ 的几何量子化 \src{371--372}。于是“带标记点的平坦联络模空间”的量子化恰好给出共形块空间 \src{373}。
\end{enumerate}

\textbf{它证明了什么、没证明什么}：
\begin{itemize}[leftmargin=1.6em]
\item 证明了（或论证到物理标准）：作用量不含度规、$k$ 的整性、$\cH_\Sigma$ 有限维、$S^3$ 上配分函数复现 Jones 多项式及其两变量推广 \src{355--356}；
\item 引用了别人的定理：Atiyah--Bott 的辛商、Borel--Weil--Bott、Segal 的共形块函子；
\item 没有证明：$\cH_\Sigma$ 与共形块空间的\textbf{函子性}（§5 只写 “It is obvious at this point” \src{397}）；半形式修正导致的能级平移 $k\to k+h^\vee$（本文略去）；Verlinde 公式的严格数学证明（本文给物理推导，数学证明 1994 年才有）。
\end{itemize}

\textbf{为什么值得读}：它示范了一种今天已成标准的做法——\emph{把拓扑不变量解释为某个量子场论的配分函数}。这件事后来发展成拓扑量子场论、Reshetikhin--Turaev 不变量、以及“用物理猜定理、用数学证定理”的整个工作方式。

\section{问题从哪来：Atiyah 的两问与 Jones 多项式的缺口}

\subsection{一个 1988 年的追问}

本文开篇引的是 Atiyah 在“去年 Hermann Weyl 讨论会”上提的两个问题 \src{351}。原话只能从 Witten 的转述读到（他说 Atiyah “might roughly be described” 这样问），原文文献 [1] 本身不在引用范围内可查。两问是：

\begin{enumerate}[leftmargin=2em]
\item 能否给 Donaldson 的四维流形理论一个物理诠释？
\item 能否给 Jones 多项式一个内蕴的三维定义？
\end{enumerate}

第一问在同一年由 Witten 自己的另一篇文章解决——用扭曲的 $N=2$ 超对称 Yang--Mills 理论 \src{351}，即参考文献 [4]（\emph{Topological quantum field theory}, CMP 117 (1988) 353）。本文处理第二问。

\begin{keypoint}
\textbf{注意}：中译本每节开头的导语框里有一句“本文解决第二个问题”。这句话是\textbf{译者所加}，不在原刊引言里。原文在引言与摘要中只说这篇理论 “gives a natural framework” 给 Jones 多项式 \src{352}，用词比“解决”弱。引用时必须回到正文。
\end{keypoint}

\subsection{Jones 多项式缺什么}

Jones 在 1983--1987 年间给出了一个纽结不变量，它有几种等价的、非常漂亮的\textbf{二维}定义：穿过投影图的切片（skein）关系、辫群表示、von Neumann 代数迹 \src{352}（文献 [5]）。但它没有\textbf{内蕴三维}的定义——不知道它在三维空间里“是什么”。

Witten 用的类比很精确 \src{352}：1930 年代人们知道量子电动力学是相对论性的，却写不出一个明显相对论性的表述；直到有了量子场论的语言，这件事才“显然”。Jones 多项式当时就处在这个位置：二维上算得出来，但不知道三维里的对象是什么。

\subsection{物理侧的动机：二维来源太多}

同期物理学家发现纽结多项式从好几个地方冒出来 \src{352}：
\begin{itemize}[leftmargin=1.6em]
\item 可解格点模型（统计力学，文献 [7] Kauffman）；
\item Yang--Baxter 方程与量子群（文献 [8] Turaev，[35] Drinfeld）；
\item 二维共形场论的单值性（文献 [11] Tsuchiya--Kanie）。
\end{itemize}

“来源太多”本身是个问题：它们之间是什么关系？Witten 的原话是，要“为这种多样性带来秩序、找到统一主题” \src{352}。更具体地，他问：\textbf{二维共形场论里，究竟是什么东西是三维的？}这个问题后来被证明是整篇文章的种子：三维的 Chern--Simons 理论沿曲面切开给出的 Hilbert 空间，正是二维共形场论的共形块。

\subsection{“拓扑不变性 = 广义协变”}

引言里有一个容易被跳过的观念转折 \src{353}：通常说一个理论“拓扑不变”，指的是配分函数不依赖度规；而广义相对论里“对度规积分”只是实现这种不变性的一种方式。Chern--Simons 理论提供了另一种：\textbf{作用量里根本没有度规}，所以不需要对度规积分，量子化出来的量自动不依赖度规。

这一步是整篇文章的立足点：如果作用量含度规，量子化后一般会得到依赖度规的量；不含度规，才有资格谈拓扑不变量。

\subsection{本文的日期与身份}

文末记：1988 年 9 月 12 日收稿，9 月 27 日修改；由 A. Jaffe 送审。p.351 的脚注标明本文是 1988 年 7 月 Swansea 国际数学物理大会报告的扩充版。所以它的“学术日历”是：先有报告，再有这篇论文；文中的一些说法（例如“量子群”评论）带着报告式的前瞻语气，不宜当定理读。

\subsection{摘要里的四个动词}

摘要 \src{351} 说：2+1 维量子 Yang--Mills 在作用量纯为 Chern--Simons 项时可以精确求解（“can be exactly soluble”）；这给 Jones 多项式一个自然框架（“gives”）；Jones 多项式可以推广到任意三维流形（“can be generalized”）；并且是可由手术计算的（“computable”）；还给 1+1 维共形场论带来“出人意料的新启示”（“shed new light”）。四个动词里没有 “prove”。这提示了全文的性质：\textbf{构造 + 论证 + 计算方案}，不是一套公理化的证明。

\begin{boundary}
\begin{itemize}[leftmargin=1.4em]
\item 已确证：Atiyah 两问的转述文本、第一问由 [4] 解决、Jones 缺三维定义、二维来源的三条线索、Swansea 版本、收稿日期。
\item 我们补的：把“二维来源过多”解释为“统一动机”，原文只说“带来秩序”。
\item 未证/待核：Atiyah 原报告的会议年份与文集信息（原文只写“去年”）；“本文解决第二问”一语不见于原刊。
\end{itemize}
\end{boundary}

\section{Chern--Simons 作用量与 $k$ 的整性}
\label{sec:cs}

\subsection{为什么不能直接用 Yang--Mills 作用量}

设 $M$ 是三维流形，$G$ 是紧单 Lie 群，$A=A_i\,\dd x^i$ 是 $M$ 上 $G$-联络（Lie 代数取值 1-形式）。规范变换的无穷小形式是 \src{353}
\begin{equation}
A_i\;\longrightarrow\; A_i-D_i\epsilon ,\qquad
F_{ij}=[D_i,D_j]=\partial_iA_j-\partial_jA_i+[A_i,A_j].
\label{eq:gauge}
\end{equation}

标准的 Yang--Mills 作用量是
\begin{equation}
L_0=\int_M\sqrt{g}\,g^{ik}g^{jl}\,\Tr(F_{ij}F_{kl}) .
\label{eq:ym}
\end{equation}
它显含度规 $g$，通过 $\sqrt g$ 与 $g^{ik}g^{jl}$ 两处 \src{353}。所以量子化后一般得不到不依赖度规的量。注意两个因子的性质不同：$\sqrt g\,\dd^3x$ 是体积元，换度规只改标度；但 $g^{ik}g^{jl}$ 把两个 $F$ 的指标缩并，换度规会改变不同分量之间的混合方式。前者可以靠“不积分体积”绕开，后者绕不开。

\subsection{Chern--Simons 作用量}

Witten 用的是 \src{354}（(1.3)）：
\begin{equation}
\begin{aligned}
L&=\frac{k}{4\pi}\int_M \Tr\Bigl(A\wedge \dd A+\frac23 A\wedge A\wedge A\Bigr)\\
 &=\frac{k}{8\pi}\int \dd^3x\,\epsilon^{ijk}\Tr\Bigl(A_i\partial_jA_k+\frac23 A_i[A_j,A_k]\Bigr).
\end{aligned}
\label{eq:cs}
\end{equation}

它的三个性质决定了一切：
\begin{enumerate}[leftmargin=2em]
\item \textbf{只依赖定向}：被积式是 3-形式，$M$ 的定向决定 $\epsilon^{ijk}$ 的符号；整个式子没有 $g$。
\item \textbf{变分给出曲率}：$\delta\int\Tr(A\dd A+\tfrac23A^3)=2\int\Tr(\delta A\wedge F)$，所以 \ours
\begin{equation}
\delta L=\frac{k}{2\pi}\int_M \Tr(\delta A\wedge F).
\label{eq:var}
\end{equation}
驻点条件是 $F=0$——平坦联络。这正是第 \ref{sec:canonical} 节相空间是“平坦联络模空间”的来源。
\item \textbf{大规范变换下不是不变，而是平移 $2\pi$ 的整数倍}：这是 $k$ 必须整数的根源，见下。
\end{enumerate}

\subsection{$k$ 的整性：大规范变换}

$\pi_3(G)\cong\ZZ$，所以规范变换 $g:M\to G$ 有绕数 $m$。在 $g$ 下作用量平移 \src{354}
\begin{equation}
L\;\longrightarrow\; L+\text{const}\cdot m .
\end{equation}
原文只写“常数乘以 $m$”，没有写出常数。我们把它算出来 \ours：对 $SU(N)$ 取 $N$ 维表示的迹，
\begin{equation}
\Delta L=\frac{k}{4\pi}\cdot\frac13\int \Tr(g^{-1}\dd g)^3
=\frac{k}{4\pi}\cdot\frac13\cdot 24\pi^2 m=2\pi k\,m ,
\label{eq:bigauge}
\end{equation}
其中用到 $\frac1{24\pi^2}\int\Tr(g^{-1}\dd g)^3\in\ZZ$。

量子理论要求 $\exp(\ii L)$ 单值，即 $\Delta L\in 2\pi\ZZ$，于是
\begin{equation}
k\in\ZZ .
\end{equation}
原文的说法是：这个条件“必要且充分” \src{355}。必要是上面这条；充分指的是：只要 $k$ 整，大规范变换就不改变 $\exp(\ii L)$。这是本文对“为什么耦合常数会量子化”的第一次回答——与二维仿射李代数中心荷的整性同源。

\begin{keypoint}
\textbf{约定陷阱}：$k$ 的整性与 $\Tr$ 的归一化绑定。若改换表示或改换 Killing 型的归一化（例如把长根平方取为 2），$k$ 的整性条件会跟着变。河野书用的是 $\mathrm{CS}=\frac1{8\pi^2}\int\Tr(A\dd A+\frac23A^3)$、配分函数 $\int e^{2\pi\ii k\,\mathrm{CS}}$ \src{河野 p.74}，与本文的 $L$ 相差一个常数因子；比较两书的 $k$ 时必须先对齐归一化。
\end{keypoint}

\subsection{Wilson 线与配分函数}

Wilson 线是沿闭曲线 $C$ 的路径序指数 \src{355}（(1.5)）：
\begin{equation}
W_R(C)=\Tr_R\,P\exp\int_C A_i\,\dd x^i .
\label{eq:wilson}
\end{equation}
它不需要度规：$A_i\dd x^i$ 是 1-形式沿曲线的积分，$\Tr_R$ 是表示空间的迹。

链环 $C_1,\dots,C_n$（带表示 $R_1,\dots,R_n$）的配分函数定义为 \src{355}（(1.6)）
\begin{equation}
Z(M;C_1,\dots,C_n)=\int \mathcal{D}A\;\; e^{\ii L}\prod_{i=1}^n W_{R_i}(C_i).
\label{eq:partition}
\end{equation}

本文的核心断言是：取 $M=S^3$、$G=SU(N)$、所有 $R_i$ 都取定义表示（$N$ 维），这个配分函数就是 Jones 多项式及其两变量推广，两个变量是 $N$ 与 $k$（解析延拓为复值）\src{355--356}；取 $G=SO(N)$、$R$ 为 $N$ 维表示则给出 Kauffman 多项式。

\textbf{这里要小心三件事}：
\begin{itemize}[leftmargin=1.6em]
\item 原文用的是 “we will claim” \src{355}，是断言，不是在这一节证明的；证明分散在 §2（弱耦合展开）与 §4（手术计算）。
\item \textbf{弱耦合极限是大 $k$}：$k\to\infty$ 时 $e^{\ii L}$ 剧烈振荡，鞍点法可用。这与 Yang--Mills 的 $1/g^2$ 展开方向相反，不能混为一谈。
\item 两变量里的 $k$ 在最终公式中要换成 $k+N$（或 $k+h^\vee$）——半形式修正，本文 §3 略去，见第 \ref{sec:canonical} 节。
\end{itemize}

\begin{boundary}
\begin{itemize}[leftmargin=1.4em]
\item 原文给出：(1.1)--(1.6)、$k$ 整性的必要性、$S^3$ 上复现 Jones 的断言。
\item 我们补的：变分公式 \eqref{eq:var} 的显式系数、$\Delta L=2\pi k m$ 的显式计算（原文只写 const$\cdot m$）。
\item 未证：$k$ 整性的“充分性”依赖测度与归一化的约定，原文未逐条论证；$S^3$ 上配分函数等于 Jones 多项式的严格证明不在本节。
\end{itemize}
\end{boundary}

\section{正则量子化：相空间、预量子化与极化}
\label{sec:canonical}

第 \ref{sec:cs} 节只写了作用量。这一节回答：为什么“沿曲面切开”能得到一个有限维的量子 Hilbert 空间？它由哪几步构成？每一步的数学前提是什么？

\subsection{第一步：切开流形，读出辛形式}

把三维流形写成 $M=\Sigma\times\mathbb{R}$，取 $A_0=0$ 的规范。作用量 \eqref{eq:cs} 中对时间是一阶的，三阶项因只剩两个空间指标而消失 \src{367}：
\begin{equation}
L=\frac{k}{8\pi}\int \dd t\int_\Sigma \epsilon^{ij}\Tr\Bigl(A_i\frac{\partial}{\partial t}A_j\Bigr).
\label{eq:sympl}
\end{equation}
这已经是“$p\dot q$”的形式，所以 $A_i$ 与 $A_j$ 是共轭变量，泊松括号直接读出 \src{367}：
\begin{equation}
\{A_i^a(x),A_j^b(y)\}=\frac{4\pi}{k}\,\epsilon_{ij}\,\delta^{ab}\,\delta^2(x-y).
\label{eq:poisson}
\end{equation}

\subsection{第二步：高斯定律与辛商}

$A_0$ 是拉格朗日乘子，变分给出高斯定律约束 \src{367}：
\begin{equation}
\epsilon^{ij}F^a_{ij}=0\quad\Longleftrightarrow\quad F=0 .
\label{eq:gauss}
\end{equation}
二维曲面只有一个独立分量，所以约束就是平坦性。它是第一类约束，生成规范变换。

经典上分两步施加：先限制到 $F=0$（平坦联络），再按规范群取商。这一步就是 Marsden--Weinstein 辛商，用的是 Atiyah--Bott 的框架 \src{367}（文献 [33]）。结果是
\begin{equation}
\cM=\Hom(\pi_1\Sigma,G)/G ,
\end{equation}
即“$\Sigma$ 上平坦联络模去规范变换”。

\begin{keypoint}
$\cM$ 是\textbf{有限维}的：$g>1$ 时 $\dim_\mathbb{R}\cM=(2g-2)d$，$d=\dim G$。它紧、带奇点、辛体积有限。半经典地看，相空间每单位体积容纳一个量子态，所以量子 Hilbert 空间有限维 \src{368}。这正是本文“可精确求解”的技术前提——如果相空间是无限维的，后面所有有限维论证都不成立。
\end{keypoint}

\subsection{第三步：预量子化}

几何量子化的标准套路是：先预量子化，再选极化。

\textbf{预量子化的前提是整性}：$[\omega]\in H^2(\cM,\ZZ)$，否则不存在预量子化线丛。Witten 在 §3.1 里\textbf{没有讨论这一条}——他直接取 $\cM_J$ 上的行列式线丛 $\cL$（带 Quillen 度规，文献 [34]），并写 $[\omega]=c_1(\cL^{\otimes k})$ \src{368}。这不是疏漏，而是构造性绕过 \ours：行列式线丛由几何本身给出，它的第一陈类自动是整的，所以“整性条件”在这个具体问题里是免费的。真正需要整性论证的是一般辛流形，不是这里。

于是预量子化 Hilbert 空间是 $\cL^{\otimes k}$ 的截面空间；但截面空间太大（同时依赖位置与动量），必须再选极化。

\subsection{第四步：极化与全纯截面}

紧相空间上不能用 $q,p$ 分离的实极化（$p_j$ 无界），所以 Witten 取 \textbf{K\"ahler 极化}：选 $\cM$ 上的复结构 $J$，把 $\cM_J$ 看成 K\"ahler 簇，量子态空间取全纯截面 \src{368}：
\begin{equation}
\cH_\Sigma^{(J)}=H^0\bigl(\cM_J,\cL^{\otimes k}\bigr).
\label{eq:hilbert}
\end{equation}
这是有限维的。

\begin{keypoint}
\textbf{极化在数学上做了什么}：它把“函数”换成“沿极化方向协变常数的截面”。实极化下这是沿拉格朗日叶的常截面；K\"ahler 极化下这是全纯截面。李思第 5.4 节把这一步讲得比本文细：先证明 $\cL$ 上的联络与辛形式相容（$[\nabla_X,\nabla_Y]=\nabla_{[X,Y]}-\ii\hbar\,\omega(X,Y)$），再定义极化，再取极化截面，最后才得到量子态空间。本文跳过了相容性这一层，直接使用结论。
\end{keypoint}

\subsection{第五步：$J$ 变时发生了什么}

$\cH_\Sigma^{(J)}$ 全纯依赖于 $J$。Witten 的断言是：这些空间拼成复结构模空间上的一个\textbf{射影平坦}向量丛 \src{369--371}；“与 $J$ 无关”的准确含义是存在典范的平坦联络，而由于共形反常，平坦只能到射影意义 \src{370}。这个联络就是 Hitchin 联络（带标记点时是 KZ 联络）；Segal 后来严格证明了它的平坦性 \src{370}（文献 [16]）。

单值性有物理意义：绕模空间一周，态空间上的变换是映射类群的作用 \src{371}。这正是二维共形场论里“模函子”的雏形。

\subsection{第六步：被略去的半形式修正}

几何量子化一般还需要半形式修正（取 $K^{1/2}$，$K$ 是典范线丛）。本文 §3 略去了它，只用 $\cL^{\otimes k}$ \src{369}。

后果是能级平移：量子化给出的组合是 $k+h^\vee$（$h^\vee$ 为对偶 Coxeter 数；$SU(N)$ 时 $h^\vee=N$），最终配分函数里的量子参数是 \src{382}
\begin{equation}
q=\exp\frac{2\pi\ii}{k+N}\qquad(SU(N)),\qquad
q=\exp\frac{2\pi\ii}{k+2}\qquad(SU(2)).
\label{eq:q}
\end{equation}
读本文时若直接用 $q=e^{2\pi\ii/k}$，会与 §4 的最终公式差一个平移。

\begin{keypoint}
\textbf{两处必须分清}：本文的平移其实是在 §2 以单圈/$\eta$ 修正的形式出现的——(2.16)(2.17) 把耦合写成 $k+c_2(G)/2$，Witten 自己指出这与二维流代数里“$k\to k+c_2/2$”是同一件事 \src{359}。半形式修正只是几何量子化语言里对同一现象的另一种描述。若按标准 $K^{1/2}$ 去推，$K_{\cM}=\mathrm{Det}^{-2h^\vee}$，给出的可能是 $k-h^\vee$——符号取决于 $\omega$ 与 $\mathrm{Det}$ 的取向约定，本文没有给。所以“$k\to k+h^\vee$”应理解为“有效组合是 $k+h^\vee$”，不要当作一条可独立推导的定理。
\end{keypoint}

\begin{boundary}
\begin{itemize}[leftmargin=1.4em]
\item 原文给出：(3.1)--(3.3)、辛商、相空间有限维、预量子化线丛、K\"ahler 极化、射影平坦性与单值性。
\item 我们补的：把“整性免费”解释为行列式线丛的构造性绕过（S2-A 的判断）；把极化步骤与李思 5.4 节对齐。
\item 未证：半形式修正导致 $k\to k+h^\vee$ 的独立推导（本文未给，三本教材也都没有讲 $K^{1/2}$ 机制，只给等价陈述）；射影平坦性的严格证明归 Segal。
\end{itemize}
\end{boundary}

\section{Wilson 线、穿刺点与 Borel--Weil--Bott}
\label{sec:wilson}

前两节只处理了“没有纽结”的理论。要让纽结进来，就得在三维流形里放 Wilson 线。这一节回答三个问题：Wilson 线怎么变成曲面上一个穿刺点？穿刺点上的约束为什么不能经典求解？Witten 用什么办法把它“反量子化”回去？

\subsection{纽结变成一个穿刺点}

把 $M$ 沿曲面 $\Sigma$ 切开时，一条穿过 $\Sigma$ 的 Wilson 线会留下一个交点 $P$。在交点附近，规范场的行为由所在表示 $R$ 决定：横向分量与 $R$ 的李代数生成元耦合。所以“带 Wilson 线的三维理论”等价于“带标记点的曲面理论”，每个标记点挂一个表示 \src{371}。

标记点上的高斯定律不再是 $F=0$，而是带源的形式 \src{371}（(3.4)）：
\begin{equation}
\frac{k}{8\pi}\,\epsilon^{ij}F^a_{ij}(x)=\sum_s \delta^2(x-P_s)\,T^a_{(s)} ,
\label{eq:markedgauss}
\end{equation}
其中 $T^a_{(s)}$ 是第 $s$ 个标记点上表示 $R_s$ 的生成元。

\subsection{为什么不能经典求解}

约束右端是\textbf{算符值}：$T^a_{(s)}$ 之间不对易，它作用在表示空间上。所以 (3.4) 不是一个可以“解出某个经典联络”的方程，而是一个量子约束 \src{372}。原文的说法是：若硬把 $A$ 当作经典对象去解，得到的不是普通联络，而是“量子群”的某种东西；但 Witten 明说这条路只是提出来，没有构造 \src{372}。

第二个困难是奇异性：点源使 $A$ 在 $P_s$ 处奇异，平坦性只在 $\Sigma\setminus\{P_s\}$ 上成立。所以相空间不是“$\Sigma$ 上平坦联络模空间”，而是“带奇点的平坦联络模空间”。

\subsection{Witten 的办法：把表示也量子化}

Witten 的解法不是解约束，而是换一个角度看表示：\textbf{把 $R$ 本身看成某个辛流形的几何量子化} \src{371--372}。

具体地，取旗流形 $G/T$（$T$ 为极大环面）。它有自然的 K\"ahler 结构；对支配整权 $\lambda$，取辛形式 $\omega_\lambda$，其 de Rham 类就是 $\lambda$。Borel--Weil--Bott 定理说
\begin{equation}
H^0(G/T,\cL_\lambda)=R_\lambda ,
\end{equation}
即 $\omega_\lambda$ 的几何量子化恰好给出最高权为 $\lambda$ 的不可约表示。注意这里 $\lambda$ \textbf{不加 $\rho$}：$\rho$ 的位移属于半形式修正那一层（$K_{G/T}^{1/2}=\cL_{-\rho}$），不在指标里。

于是可以构造一个更大的相空间 $M_{P_i,R_i}$：在 $\Sigma\setminus\{P_i\}$ 上取平坦 $G$ 联络，在每个 $P_i$ 处把结构群约化到 $T$，参数由 $\lambda_i$ 给出。这个空间上的经典约束 (3.4) 有意义；量子化之后，得到的恰好是 Segal 的流代数共形块空间 \src{372--373}。

\begin{keypoint}
\textbf{这一步的逻辑}：不是在三维理论里“解”穿刺约束，而是把“表示”从外部数据变成相空间的一部分。原来 Wilson 线 $W_R(C)=\Tr_R P\exp\int_C A$ 里的 $\Tr_R$ 与 $R$，现在由 $G/T$ 这一层的量子化提供。Witten 在 p.371 说这是 “should” 的做法，并承认 “deferring a fuller treatment”——所以这是\textbf{纲领 + 部分实现}，不是完整构造。
\end{keypoint}

\subsection{可积表示条件}

只有“可积”的表示才有非零截面：$SU(2)$ 时 $2j\le k$，一般地 $\lambda$ 属于 level $k$ 的仿射 Weyl 群基本域。超出这个范围的表示在共形块里“退耦”（对应维数为零）\src{372}。

\textbf{这个条件不是 Borel--Weil--Bott 的推论}——BWB 对任意支配整权都成立。它来自仿射李代数：level $k$ 的可积表示要求 $(\lambda,\theta)\le k$（河野命题 1.5，书 p.11）。原文把它写在 p.372 的括号里，一笔带过；Moore 中译 p.120 与上野 p.39 才明写。

\subsection{记号对照}

读原文时容易把几个“$\lambda$”混起来：
\begin{center}
\small
\begin{tabular}{@{}lll@{}}
\toprule
对象 & 记号 & 说明 \\
\midrule
最高权 & $\lambda$ & 表示 $R_\lambda$ 的最高权；BWB 里不加 $\rho$ \\
辛形式 & $\omega_\lambda$ & $G/T$ 上，$[\omega_\lambda]=\lambda$ \\
共形块 & $H^0(\cM_J,\cL^{\otimes k})$ & 带标记点时是子空间，由 $\lambda_i$ 决定 \\
仿射 level & $k+h^\vee$ & 半形式修正后；本文多处写 “level $k$” 是简写 \\
\bottomrule
\end{tabular}
\end{center}

\begin{boundary}
\begin{itemize}[leftmargin=1.4em]
\item 原文给出：(3.4)、BWB 实现、扩大相空间、引 Segal、可积性括号。
\item 我们补的：把“$\lambda$ 不加 $\rho$”与半形式修正分开（$\rho$ 属于 $K^{1/2}$）；把 $2j\le k$ 归因于仿射代数而非 BWB。
\item 未证：把 Wilson 线换成“$S\to G/T$ 的映射的泛函积分”这一纲领，原文只有 “should” 与 “deferring a fuller treatment”，没有构造；“量子群”评论同样没有构造。
\end{itemize}
\end{boundary}

\section{共形块与 WZW：这个对应是什么级别}
\label{sec:blocks}

这是全篇最容易被说过头的地方。用户会问：Witten 到底证明了“三维 Chern--Simons 理论 = 二维 WZW 模型”，还是只证明了“两边的态空间维数一样”？答案在两者之间，而且必须把“对象”“函子”“范畴”三个层级分开。

\subsection{两条独立的路线}

WZW 一侧：一个二维共形场论，由 Wess--Zumino 项加 WZW 项定义，其手征代数（仿射李代数）的表示给出共形块。BPZ 方程、Tsuchiya--Kanie 的工作、Segal 的模函子，是这一侧的数学语言 \src{352}（文献 [11][16][32][39]）。

CS 一侧：第 \ref{sec:canonical} 节把三维理论沿 $\Sigma$ 切开，得到 $\cH_\Sigma^{(J)}=H^0(\cM_J,\cL^{\otimes k})$。

问题是：这两边是不是同一个东西？

\subsection{原文说了什么}

Witten 在 §3.1 说这些态空间“与 Segal 的模函子是同一物” \src{366}；在 §3.2 直接把 $\cH_\Sigma$ 与 BPZ 的共形块空间等同 \src{370}。英文用的是 “\textbf{are}”，不是 “isomorphic to”——从措辞看，他认为是同一个对象，不是碰巧同构。

但证据只到\textbf{对象级}：两边都等于 $H^0(\cM_J,\cL^{\otimes k})$，带同一个射影平坦丛 \src{369--370}。他没有证明函子性（缝合相容），§5 只有一句 “It is obvious at this point” \src{397}。

原文最强的\textbf{数学}断言在 §5 \src{396}：按 Segal 的看法，把这个辛结构量子化 $LG/G$，得到的正是圈群的基本不可约最高权表示。但这一句归功于 Segal（文献 [39]），而且限于 genus 0 与基本表示。

\subsection{三个层级}

\begin{center}
\small
\begin{tabular}{@{}p{2.3cm}p{5.2cm}p{6.2cm}@{}}
\toprule
层级 & 断言 & 原文证据 \\
\midrule
对象级 & 两边的态空间是同一个向量空间 & p.370 “are”；两边都是 $H^0(\cM_J,\cL^{\otimes k})$ \\
函子级 & 复结构变化时，两边给出同一个射影平坦丛；缝合相容 & p.370 只给丛与联络；缝合只到 p.397 “It is obvious” \\
范畴级 & 两个理论作为 TQFT/CFT 等价 & 未证；framing、边界条件、配分函数层面都有余量 \\
\bottomrule
\end{tabular}
\end{center}

所以正确的说法是：\textbf{强于“维数相同”，弱于“范畴等价”}。写“Witten 证明了两个理论等价”是错的；写“只是维数巧合”也错。Moore 中译的措辞更精确：向量空间同构，在 K\"ahler 量子化时是模空间上全纯向量丛的等同 \src{Moore 中译 p.76}。

\subsection{level 的账：裸 $k$ 与有效 $k+h^\vee$}

WZW 侧的一切公式里出现的是 $k+h^\vee$，不是裸 $k$。来源不是“两个理论能级不同”，而是 \S2 的单圈修正 \src{359}：(2.16)(2.17) 里 $\eta$ 的度规依赖把耦合平移为 $k+c_2(G)/2$，而 $c_2(SU(N))=2N$，故平移量是 $h^\vee=N$。Witten 自己说，这“似乎解释了下述事实：1+1 维流代数的许多公式中，量子修正的效果正是把 $k$ 替换为 $k+c_2(G)/2$” \src{359}。

在几何量子化的语言里，同一个平移来自半形式修正（第 \ref{sec:canonical} 节末）。最终 $SU(N)$ 的量子参数是 \eqref{eq:q}，$SU(2)$ 是 $q=e^{2\pi\ii/(k+2)}$ \src{382}。

\begin{keypoint}
\textbf{术语建议}：把 $k$ 叫“裸能级”、$k+h^\vee$ 叫“有效能级”。本地译本有一处写“WZW 能级是 $k+N$”，按这个口径是“有效能级”；若严格按模 $S$ 矩阵的 $q$ 定义“能级”，则 $k+h^\vee$ 才配称能级。两种说法不冲突，但混用会与可积性判据 $2j\le k$ 打架（后者用的是裸 $k$）。
\end{keypoint}

\subsection{KZ 方程：正文没有，但就在旁边}

KZ（Knizhnik--Zamolodchikov）方程在本文正文里\textbf{零出现}；最接近的是 p.370 引 BPZ 的微分方程 \src{370}。但教材里 KZ 正是 §3.2 那条射影平坦联络在 $g=0$、$n$ 个标记点时的显式形式：共形块丛 $\mathcal{E}_\lambda\to\mathrm{Conf}_n(\mathbb{CP}^1)$ 上的联络 $\nabla=\dd-\sum_i L^{(i)}_{-1}\dd z_i$，水平截面满足 KZ 方程，系数是 $\frac{1}{k+2}\frac{\Omega^{(ij)}}{z_i-z_j}$ \src{河野 p.25--28}。

所以本文与 KZ 的关系是：本文给“存在一条平坦联络”，KZ 给“这条联络长什么样”。不要把后者读进本文。

\begin{boundary}
\begin{itemize}[leftmargin=1.4em]
\item 原文给出：对象级等同（p.370）、射影平坦丛（p.370）、Segal 定理的引用（p.396）、三点维数 $N_{ijk}$（p.373--374）。
\item 我们补的：三层级划分；“裸能级/有效能级”的术语建议；KZ 与本文关系的定位。
\item 未证：函子性（p.397 只有 “obvious”）；范畴级等价；半形式修正的符号（$K^{1/2}$ 给 $k-h^\vee$ 还是 $k+h^\vee$，取决于 $\omega$ 与 $\mathrm{Det}$ 的取向约定，原文未给）。
\end{itemize}
\end{boundary}

\section{Verlinde 公式：从共形块到融合规则}
\label{sec:verlinde}

用户原话：“这个公式怎么从共形块推出来？它的物理意义是什么？”

这一节先写公式是什么，再写 Witten 的推导链，最后写它的地位——因为这条公式是本文里“物理推导”与“数学定理”距离最近、也最容易被混淆的地方。

\subsection{公式}

对可积表示 $i,j,k$（$SU(2)$ 时 $2j\le k$），融合系数
\begin{equation}
N_{ijk}=\sum_l \frac{S_{il}S_{jl}S_{kl}}{S_{0l}} ,
\label{eq:verlinde}
\end{equation}
其中 $S$ 是模变换矩阵，由 $q=\exp\frac{2\pi\ii}{k+N}$ 决定 \src{382}。$N_{ijk}$ 有三重身份：
\begin{enumerate}[leftmargin=2em]
\item 融合规则 $R_i\otimes R_j=\bigoplus_k N_{ij}{}^k R_k$ 里的重数；
\item $S^2$ 上三个标记点（表示 $R_i,R_j,R_k$）的 Hilbert 空间维数，$\dim\cH_{S^2}(P_i,P_j,P_k)=N_{ijk}$ \src{373--374}；
\item 三点函数空间的维数 $\dim\Inv(R_i\otimes R_j\otimes R_k)$。
\end{enumerate}

\subsection{推导链：十个步骤，三处输入}

Witten 在 §4.5 “Some Concrete Surgeries”（起 p.388）里用一段斜体小标题 “Proof of Verlinde's Conjecture” 展开推导 \src{391}。注意它是\textbf{段首小标题}，不是节标题——这也提示读者：这一段是本文对 Verlinde 猜想的“物理重证”，而猜想本身已由 Moore--Seiberg 证明 \src{391}。推导链可以压缩成十步，但\textbf{真正的输入只有三处} \ours：

\begin{enumerate}[leftmargin=2em]
\item \textbf{手术律} (4.27) \src{385}／(4.28) \src{386}：沿链环 $C$ 做手术时，配分函数按 $S$ 矩阵展开。这是链的动力学输入。
\item \textbf{$N_{abc}$ 的定义} (4.36) \src{388}：把 $N_{abc}$ 定义为 $S^2\times S^1$ 上三标记点的配分函数，也就是三维 Hilbert 空间的维数。这是“把要求解的量翻译成可算的量”。
\item \textbf{秩一性} (4.49) \src{393}：$\dim\cH=1$ 时的归一化。这是最后一步把比值定下来的条件。
\end{enumerate}

中间的关键一步是第 6 步：把“平行线”变成“链接线”，即把 $\delta_{ij}$ 换成 $S_{ij}$——这一步用的是广义手术公式 (4.28)，也是全链唯一需要它的地方 \ours。

结果是 \src{393}（(4.51)）：
\begin{equation}
\sum_k S_i{}^k\,N_{kl}{}^m=\frac{S_{il}S_i{}^m}{S_{0i}} .
\label{eq:verlinde-key}
\end{equation}

\begin{keypoint}
\textbf{这条推导链没有用到共形块模型的任何动力学}。它只用“配分函数在手术下如何变换”。所以 (4.51) 的内容完全可以从三维理论的自洽性推出——这正是 Verlinde 公式“物理意义”的核心：\textbf{融合规则不是额外输入，它被模变换矩阵 $S$ 完全决定}。
\end{keypoint}

\subsection{两种读法}

(4.51) 可以朝两个方向读 \ours：

\begin{itemize}[leftmargin=1.6em]
\item \textbf{读法 (a)：融合规则 $\to S$}。若已知 $N_{ij}{}^k$，则 (4.51) 说 $S$ 把它们对角化。这是“陈述”。
\item \textbf{读法 (b)：$S\to$ 幂等元}。定义 $w_i=S_{0i}\sum_m S_i{}^m v_m$，则 (4.51) 等价于
\begin{equation}
w_i w_j=\delta_{ij}w_j ,
\end{equation}
即 $w_i$ 是 Verlinde 代数 $v_i v_j=\sum_k N_{ij}{}^k v_k$ 的一组正交幂等元，$N_{ijk}$ 由 $S$ 显式给出。这是“构造”，也是 Moore--Seiberg 之后的通行读法 \src{393}。
\end{itemize}

两种读法数学等价，但解释不同：(a) 把 $S$ 当作派生量，(b) 把 $S$ 当作基本量、融合规则是它的影子。物理上 (b) 更自然——$S$ 来自模变换，是几何的；融合规则是代数后果。

\subsection{物理意义}

一句话：\textbf{在三维 Chern--Simons 理论里，二维融合规则被三维的模变换完全决定}。所以：

\begin{itemize}[leftmargin=1.6em]
\item Verlinde 公式不是“又一个公式”，而是“三维理论自洽性对二维数据的约束”；
\item 它给出了 $N_{ijk}$ 的算法：只要知道 $S$，不用解任何共形场论就能算融合系数；
\item 反过来，它把 $S$ 矩阵变成主角——这也解释了为什么后续的模张量范畴、Reshetikhin--Turaev 不变量都以 $S$ 为中心。
\end{itemize}

\subsection{地位：物理推导 vs 数学定理}

原文的说法是：Verlinde 提出猜想 [12]，Moore--Seiberg 给出证明 [13]，“我们给出新证明” \src{391}。这里的“证明”是物理推导——用配分函数在手术下的形式性质，不是从 CS 测度出发的数学证明。

数学侧的四级路线要分清 \src{Schottenloher 书 p.177；上野 p.106}：
\begin{center}
\small
\begin{tabular}{@{}p{3.4cm}p{4.2cm}p{5.4cm}@{}}
\toprule
层面 & 谁 & 状态与工具 \\
\midrule
物理推导 & Witten §4.5（本文） & TQFT 公理下的推导，自洽可复算 \\
物理证明 & Moore--Seiberg 1989 & 首证；Witten 自认“新证明” \\
数学证明（几何） & TUY 1989；Faltings 1994 & 共形块丛 + KZ/Hitchin；rank-level 对偶 \\
组合/范畴 & RT 1991；模张量范畴 & 独立路线，以 $S$ 为中心 \\
现代代数证明 & 融合环 = Frobenius 代数 & 谱分解，把“自然数”变成代数事实 \\
\bottomrule
\end{tabular}
\end{center}

Schottenloher 的原话是“这些证明都不容易理解”：KZ 单值性加非交换 Hodge 理论是重机器；右端是 $\sin$ 的幂次和，“是自然数”本身都不显然 \src{Schottenloher 书 p.177}。

\begin{keypoint}
\textbf{一个常见范畴错误}：河野书的命题 1.22（书 p.37）与定理 2.16（书 p.60）用的是同一个 $S$ 矩阵，但前者是融合代数的基变换（代数侧，就是 Verlinde 公式），后者是 Dehn 手术核（几何侧，定义 Witten 不变量）。定理 2.16 的证明拿 1.22 当砖，但把手术公式当成“Verlinde 的证明”是张冠李戴 \src{河野 p.60--61}。
\end{keypoint}

\subsection{一个可当场做的验证}

$SU(2)_k$ 的 $S$ 矩阵有显式公式
\begin{equation}
S_{jl}=\sqrt{\frac{2}{k+2}}\,\sin\frac{\pi(2j+1)(2l+1)}{k+2},
\qquad j,l=0,\tfrac12,\dots,\tfrac{k}{2}.
\label{eq:smatrix}
\end{equation}
它由 $q=\exp\frac{2\pi\ii}{k+2}$ 与 (5.29)(5.30) 给出 \src{上野 p.103；河野 p.60}。

\textbf{$k=1$}：只有 $j=0,\frac12$，
\begin{equation}
S=\frac{1}{\sqrt2}\begin{pmatrix}1&1\\1&-1\end{pmatrix}.
\end{equation}
代入 \eqref{eq:verlinde} 逐个算：
\begin{align}
N_{\frac12,\frac12,0}&=\sum_l\frac{S_{\frac12 l}S_{\frac12 l}S_{0l}}{S_{0l}}
=\sum_l S_{\frac12 l}^2=\frac12+\frac12=1,\\
N_{\frac12,\frac12,\frac12}&=\sum_l\frac{S_{\frac12 l}^3}{S_{0l}}
=\frac{(1/\sqrt2)^3}{1/\sqrt2}+\frac{(-1/\sqrt2)^3}{1/\sqrt2}
=\frac12-\frac12=0 .
\end{align}
后一式正是截断的直接后果：$\frac12\otimes\frac12=0$，而 $j=1$ 的耦合被 $2j\le k$ 切掉，所以三点耦合为零。\textbf{注意}：这里“$N_{111}=0$”里的 $1$ 指 $j=1$（不可积），不要与“$j=\frac12$ 记作 1”混用。

\textbf{$k=2$}：$j=0,\frac12,1$，由 \eqref{eq:smatrix} 得 $S_{00}=\frac12$、$S_{0\frac12}=\frac{1}{\sqrt2}$、$S_{01}=\frac12$、$S_{\frac12\frac12}=0$、$S_{\frac12 1}=-\frac{1}{\sqrt2}$、$S_{11}=\frac12$。于是
\begin{equation}
N_{\frac12,\frac12,0}=1,\qquad N_{\frac12,\frac12,1}=1,
\end{equation}
即 $\frac12\otimes\frac12=0\oplus1$；同法得 $\frac12\otimes1=\frac12$、$1\otimes1=0\oplus1$。与截断融合规则逐条对上 \ours（这组数由 S5-B 用 python 独立复算过）。

\textbf{这个例子说明两件事}：一是 Verlinde 公式确实“只要 $S$ 就能算融合”；二是可积性截断必须同时作用在求和指标与结果指标上，否则会得到非整数。


\begin{boundary}
\begin{itemize}[leftmargin=1.4em]
\item 原文给出：§4.5 的十步链、(4.51)、两种读法、(4.36) 的定义、$SU(2)$ 的 $q$；段首小标题 “Proof of Verlinde's Conjecture” 在 p.391，属 §4.5（节标题 Some Concrete Surgeries，p.388）。
\item 我们补的：三处输入的归纳；两种读法的方向区分；$SU(2)_{k=1,2}$ 的数值验证；河野 1.22 与 2.16 的范畴区分。
\item 未证：数学严格证明（TUY/Faltings 只按二手转述，未读原文）；$S$ 矩阵归一化的细节；标架反常 (4.41) 在一般手术下的追踪——本文主链只用标准 $S$，Witten 明说 (4.41) 那种修正“在本文的简单应用中不必要” \src{390--391}。
\end{itemize}
\end{boundary}

\section{2d/3d 对偶：CS/WZW 与 AdS$_3$/CFT$_2$ 是两码事}
\label{sec:duality}

用户原话：“这里的对偶跟 AdS$_3$/CFT$_2$ 有什么区别？CS/WZW 对偶和全息对偶是两码事吗，这个得搞清楚。”

答案是：\textbf{是两码事}，而且不是同一种“对偶”。

\subsection{两种关系的类型不同}

\begin{center}
\small
\begin{tabular}{@{}p{1.9cm}p{5.0cm}p{5.0cm}@{}}
\toprule
维度 & CS/WZW（本文，1988） & AdS$_3$/CFT$_2$（1997 后） \\
\midrule
体理论 & 三维 Chern--Simons：无度规（作用量 (1.3) 不含 $g$）\src{354}，态空间 $\cH_\Sigma$ 有限维 \src{366--367} & 三维量子引力（含度规，渐近 AdS$_3$ 边界条件）\\
边界理论 & 二维 WZW／手征代数；其共形块空间（原文 p.370 以 “conformal blocks” 名之）\src{370} & 二维 CFT（一般大 $c$）\\
关系类型 & 量子化态空间 = 共形块空间：\textbf{态空间等同} \src{368--370}；三维侧只在 $X\times S^1$ 以 $Z=\dim\cH_X$ 接触二维数据 \src{387--388} & 体配分函数 = 边界配分函数：\textbf{动力学等价} \\
体自由度 & 无度规自由度；$\cH_\Sigma$ 有限维 \src{366--367} & 有：引力自由度、BTZ 黑洞、黑洞熵 \\
引力子 & 无（理论里没有度规） & 有 \\
中心荷 & $c=k\dim G/(k+h^\vee)$ \src{390 (4.41)}；$k\to\infty$ 时 $c\to\dim G$ \src{361}；$k\in\mathbb Z$ \src{354} & Virasoro 渐近对称，$c=3l/2G$（Brown--Henneaux），连续\\
可证性 & 可逐条检验（辛商、Narasimhan--Seshadri、BWB） & 无定理；只有弱耦合/对称性/可积性检验 \\
全息字典 & 不存在“体场↔边界算子”字典 & 存在（GKPW 型）\\
\bottomrule
\end{tabular}
\end{center}

\subsection{唯一的接触点，以及为什么它不是等同}

还要说清一件事：本文的“三维—二维接触”只有一条通道——在 $X\times S^1$ 型流形上，三维配分函数等于 $X$ 上 Hilbert 空间的维数，即 $Z(X\times S^1)=\dim\cH_X$ \src{387--388}。\textbf{没有“体配分函数 = 边界配分函数”这样的等式}；这正是它不能叫全息对偶的技术原因。

确有一个接触点：\textbf{三维引力可以写成 Chern--Simons 理论}（规范群取 $ISO(2,1)$ 或其覆盖群）。Witten 在本文只一句带过，细节在姊妹篇 \src{397}（文献 [38]，预印本 HEP-88/32）。

但这个改写\textbf{不是}等同：它要求\textbf{非紧}规范群，要求特定边界条件；本文的 CS 用\textbf{紧}群、没有边界引力自由度。改写保留的是经典运动方程，并没有把三维引力变成“有限维态空间”的拓扑场论。

\begin{keypoint}
\textbf{原文从未出现 AdS、从未出现“全息”这个词}（1988 年还没有这个说法）。本文 §5 结尾只说“正是这类考虑促成了对 2+1 维引力的研究” \src{397}——这是动机链，不是等价性陈述。
\end{keypoint}

\subsection{为什么容易混}

两个原因。第一，两者都叫“3d/2d 对应”，字面上都是“三维理论对应二维理论”。第二，CS/WZW 的边界理论确实是二维 CFT，而 AdS$_3$/CFT$_2$ 的边界也是二维 CFT——但前者的“二维理论”是\textbf{态空间的组织方式}（共形块），后者的“二维理论”是\textbf{另一个动力学理论}。

一句话区分：\textbf{CS/WZW 是同一套量子化数据的两种写法；AdS$_3$/CFT$_2$ 是两个理论之间的动力学等价。}

\begin{boundary}
\begin{itemize}[leftmargin=1.4em]
\item 已确证：本文的对应是态空间级 \src{368--369,394--397}；原文无 AdS/全息字样；姊妹篇 [38] 是 3d 引力=CS 的出处。
\item 我们补的：对照表的六个维度（体自由度、引力子、中心荷来源、可证性、字典）。
\item 未核（后世文献卷页）：Brown--Henneaux 的 $c=3l/2G$、Maldacena 1997、GKPW 字典；Achúcarro--Townsend 1986。
\end{itemize}
\end{boundary}

\section{延伸：纽结同调的范畴化}
\label{sec:categorification}

这一节超出本文范围，但用户点名要作为延伸。问题：Witten 1989 与 Khovanov 同调是什么关系？

一句话答案：\textbf{不是推论关系，但 1989 年的 CS 是“物理来源”的起点}。Khovanov 同调把 Jones 多项式从“一个数”升级成“一族双分次向量空间”，Jones 是它的分次 Euler 特征；而 Witten 在 2011 年把这件事的物理来源放到了高一维（4d/5d），并且明说“范畴化”在他的口径里就是“能从高一维理论导出”。

\subsection{范畴化：数 $\to$ 复形}

Jones 多项式是一个数（Laurent 多项式）。Khovanov 同调把每个链环 $L$ 换成一个双分次同调群 $H^{i,j}(L)$，使得（在适当归一化下；不同作者的 $q$ 与 $q+q^{-1}$ 因子约定不同）
\begin{equation}
\sum_{i,j}(-1)^i q^j\dim H^{i,j}(L)=J_L(q).
\end{equation}
即 Jones 多项式是它的 \textbf{Euler 特征} \ours。所以 Khovanov \textbf{没有证明} Witten 的 (1.6)：他假定 Jones 多项式已知，构造一个以它为 Euler 特征的同调。这与本文不在同一层。

升级换来的东西有两样，都是 1989 年没有的：
\begin{itemize}[leftmargin=1.6em]
\item \textbf{函子性}：链环协边 $\Sigma:K\to K'$ 给出线性映射 $\Phi_\Sigma:\mathrm{Kh}(K)\to\mathrm{Kh}(K')$，拼接相乘；1989 的 TQFT 只有“数 $\to$ 数”；
\item \textbf{解释整数性}：Jones 的系数为什么是整数？CS 路径积分给 Laurent 多项式，但不说这件事；Khovanov 把系数解释成分次向量空间的维数——Witten 2011 把这句写成 §1.1.2 的标题 “What Chern--Simons Theory Doesn't Explain” \src{Witten 2011 p.8}。
\end{itemize}

\subsection{物理来源是后世的，且“范畴化 = 升一维”}

本文给的是 Jones 的 3d CS 描述。Khovanov 同调的物理来源出现在 2007--2011 年：
\begin{itemize}[leftmargin=1.6em]
\item Witten 2011：D3 膜终止于带 $\theta$ 角的 NS5 膜，两支 S/T 对偶给出 Jones 与 Khovanov 的流形无关描述（4/5 维椭圆方程解的计数），arXiv:1101.3216；
\item Kapustin--Witten 2007：4d $\cN=4$ SYM 的拓扑扭曲族，提供 Wilson--'t Hooft 与面算子框架；它讲几何 Langlands，Khovanov 不是它的推论，只是同一框架的另一应用，arXiv:hep-th/0604151。
\end{itemize}

Witten 自己的点题句值得原样记住：
\begin{quote}
``the fact that the four-dimensional equations (5.38) can be `lifted' \dots\ to five dimensions is part of the reason that Chern--Simons gauge theory can be `categorified,' which is just a fancy way to say that it can be derived from a theory in one dimension higher.''\\[2pt]
\hfill --- Witten 2011, p.107
\end{quote}

这是\textbf{并列的两套物理实现}（3d CS 对 Jones；4d/5d 对 Khovanov），不是同一套的延伸。任何“Witten 1989 已经预示 Khovanov”的说法都是时代错置。

\subsection{本文有没有范畴化暗示：没有}

全文没有 homology、chain complex、categorification 一类构造，不变量始终是数 $Z(M;\text{links})$。唯一擦边处是 p.372：解带标记点约束需要“$q$ 数联络”，其和乐 “would presumably be an element of a `quantum group', not an ordinary classical group. Indeed, it seems likely that\dots” \src{372}——“presumably / seems likely” 双重 hedging，且说的是量子群，不是同调。

但本文有“半截”结构：对曲面 $\Sigma$ 给向量空间 $\cH_\Sigma$（§3），对三维流形给数——已是 Atiyah TQFT 的雏形，只是取值为向量空间而非复形，缺 $\ZZ$-分次与微分。范畴化要做的正是把 $\cH_\Sigma$ 换成复形、把 $Z(M)$ 换成上同调。

\subsection{要把本文理论范畴化，缺什么}

五条开放问题 \ours：
\begin{enumerate}[leftmargin=2em]
\item \textbf{一维提升}：需要 4d（或 5d/6d）理论，其边界约化给 CS；3d 理论自身没有“多出来的一维”放同调方向。
\item \textbf{双分次}：需要 $Q^2=0$ 与 $U(1)\times U(1)$ 给 $(i,j)$ 分次；纯 CS 只有规范约束，没有 $R$-对称性。
\item \textbf{标架/半形式反常}：CS 配分函数依赖标架（更一般地 two-framing），换标架乘 $q$ 的幂；1989 用 $\eta$ 不变量 (2.16)(2.17) 处理相位。要得到良定义的 $\ZZ$-分次复形，必须把反常典范化。
\item \textbf{表示范畴升级}：Wilson 线只带一个可积表示，对应一个模张量范畴；范畴化要求把它换成 2-范畴。一般 $G$ 未知，$A$ 型有 Khovanov--Rozansky/Soergel 路线。
\item \textbf{level 与 $q$ 的错位}：1989 的不变量定义在 $q=e^{2\pi\ii/(k+N)}$（固定 level），Khovanov 是 $q$ 的多项式（所有 level 同时）；“逐 level 范畴化”与 Khovanov 同调不是同一件事。
\end{enumerate}

T6 的真空也在这里：1989 的“CS $\to$ Jones”与辛几何路线的“$\mathrm{Kh}_{\rm symp}\to$ Khovanov”之间，目前没有一条已知定理把 CS 的 Wilson 线直接变成 Khovanov 的链复形。本地线索：\texttt{functoriality\_work/work/通读笔记.md:55}（在 \texttt{C:\textbackslash Users\textbackslash admin\textbackslash Downloads} 下）记 Seidel--Smith 引理 44、Abouzaid--Smith 2019 证明 $\mathrm{Kh}_{\rm symp}$ 与 Khovanov 同构——这是\textbf{纯数学}路线，不过 CS 路径积分。

\begin{keypoint}
\textbf{一句话}：1988 年的问题是“Jones 的三维定义是什么”；2000 年的问题是“怎样把 Jones 升级成同调”。后者更精细，但不回答前者。把两者放进同一章，只能作为“这篇论文开启的后续地图”，不能作为“这篇论文的内容”。
\end{keypoint}

\begin{boundary}
\begin{itemize}[leftmargin=1.4em]
\item 已确证：Euler 特征退回 Jones；Witten 2011 的 arXiv 号与 p.107 原句；Kapustin--Witten 的 arXiv 号；Abouzaid--Smith arXiv:1504.01230。
\item 未核：Khovanov 2000 的期刊卷页（Duke 101 (2000) 359--426 已核）；Khovanov 立方体构造的符号与 $q$ 次平移（未读原文）；“纯 CS 只有 ghost 数一维”是我们的重建。
\item 结论：本文不涉及范畴化；写成“延伸”时须整段标“后世”。
\end{itemize}
\end{boundary}

\section{边界：这篇论文没有做什么}
\label{sec:boundary}

前面八章按问题顺序走完了全文。这一章只做一件事：把“本文证明了什么”与“后世补了什么”分开列清。读经典论文最容易犯的错，是把后人补的定理算到作者头上，或者反过来把作者的断言当成已证。

\subsection{本文自己的留白}

\begin{enumerate}[leftmargin=2em]
\item \textbf{一般手术下的标架因子}。§4 对标准手术（模变换 $S$：$S^2\times S^1\to S^3$）算过相位，但对一般整系数 Dehn 手术，“如何在手术下系统地追踪 $e^{2\pi\ii c/24}$ 因子”，原文明确说留待另文 \src{391}。
\item \textbf{函子性}。态空间的对象级等同有 \src{370}，但“沿曲面粘合时共形块如何配对”只有 §5 的 “It is obvious at this point” \src{397}。
\item \textbf{半形式修正的系统处理}。本文只用 $\cL^{\otimes k}$，能级平移以 §2 的 $\eta$ 修正形式出现 \src{359}；以 $K^{1/2}$ 语言给出的独立推导没有。
\item \textbf{“量子群”}。p.372 提到把 $A$ 当经典对象硬解会得到量子群式的东西，但没有构造。
\item \textbf{一般群与一般表示}。文中反复以 $SU(N)$ 定义表示与 $SU(2)$ 为例；对一般单群、一般可积表示，只给框架。
\end{enumerate}

\subsection{后世补了什么}

\begin{center}
\small
\begin{tabular}{@{}p{4.4cm}p{4.2cm}p{4.6cm}@{}}
\toprule
问题 & 本文状态 & 后世 \\
\midrule
三维流形不变量的严格定义 & 手术算法 + 算例；一般手术标架留白 & Reshetikhin--Turaev 1991（组合定义）\\
TQFT 公理化 & 只用切割-粘贴的手工版本 & Atiyah 1988 公理；后世函子化\\
共形块 = 态空间的函子化 & 对象级 + Segal 定理引用 & TUY／后世 TQFT 文献\\
Verlinde 公式的数学证明 & 物理推导（§4.5） & TUY 1989 → Faltings 1994\\
Khovanov 同调 & 不涉及 & Khovanov 2000；物理来源 Witten 2011\\
2+1 维引力 & 姊妹篇 [38] 一句带过 & 与 Brown--Henneaux 结合成 AdS$_3$ 图像（1990s）\\
\bottomrule
\end{tabular}
\end{center}

\subsection{怎样读这篇论文：三遍法}

\begin{enumerate}[leftmargin=2em]
\item \textbf{第一遍（半天）}：读引言、§1、§3 的节首与节末、§5。目标是能说出“作用量是什么、相空间是什么、Hilbert 空间是什么、和共形块什么关系”。
\item \textbf{第二遍（两三天）}：读 §2 与 §4 的推导主干，重点是 (2.16)(2.17) 的 $\eta$ 修正、(4.21) 的 $q$、(4.28)--(4.51) 的手术链。目标是能复算 $S^3$ 上的 Jones 多项式与 $S^2$ 上三点维数。
\item \textbf{第三遍（按需）}：回到具体公式核对约定（$k$ 的归一化、$\Tr$ 的表示、半形式修正），并区分“原文说／我们补／后世”。
\end{enumerate}

\begin{keypoint}
\textbf{一句总结}：这篇论文给的是“一个可精确求解的构造 + 一整套自洽计算 + 若干断言”。把它读成“已经严格的定理集”会处处别扭；把它读成“物理学家给出的、后来被数学严格化的纲领”，每一节的位置就都清楚了。
\end{keypoint}

\begin{boundary}
\begin{itemize}[leftmargin=1.4em]
\item 本章所有“本文状态”均有前文锚点；“后世”一栏的卷页，凡未核者已在正文标出。
\item 未核：RT 1991、TV 1992、Khovanov 2000 的期刊卷页；Atiyah TQFT 公理原文卷页。
\end{itemize}
\end{boundary}


\appendix
\section{符号表与约定}
\label{app:symbols}

\begin{center}
\small
\begin{tabular}{@{}ll@{}}
\toprule
符号 & 含义 \\
\midrule
$M$ & 三维流形（通常紧、定向、无边）\\
$\Sigma$ & 二维曲面（切开 $M$ 所得）\\
$G$ & 紧单 Lie 群；$d=\dim G$；$h^\vee$ 对偶 Coxeter 数 \\
$A$ & $G$-联络，Lie 代数取值 1-形式 \\
$F$ & 曲率 2-形式，$F=\dd A+A\wedge A$ \\
$k$ & Chern--Simons 耦合常数（整数）\\
$L$ & Chern--Simons 作用量 \eqref{eq:cs} \\
$W_R(C)$ & 沿 $C$、表示 $R$ 的 Wilson 线 \\
$\cM$ & $\Sigma$ 上平坦联络模空间 $\Hom(\pi_1\Sigma,G)/G$ \\
$\cM_J$ & 选复结构 $J$ 后的 K\"ahler 模空间 \\
$\cL$ & 行列式线丛（预量子化线丛）\\
$\cH_\Sigma$ & 沿 $\Sigma$ 切开所得的量子 Hilbert 空间 \\
$H^0(\cM_J,\cL^{\otimes k})$ & K\"ahler 极化下的量子态空间（全纯截面）\\
$R_i,\ \lambda$ & 表示与最高权；$\omega_R$ 为 $G/T$ 上的辛形式 \\
$S_{ij}$ & 模变换矩阵（Verlinde 公式中的 $S$ 矩阵）\\
$N_{ijk}$ & 融合系数（三点函数维数）\\
\bottomrule
\end{tabular}
\end{center}

\textbf{约定}：
\begin{itemize}[leftmargin=1.6em]
\item 本文的 $L$ 与河野书的 $\mathrm{CS}$ 差一个常数因子（$L=2\pi k\,\mathrm{CS}_{\text{河野}}$），比较 $k$ 前先对齐归一化。
\item 页边注“原文 p.$N$”指 CMP 印刷页码；中译本 PDF 页 $=$ 印刷页 $-350$ 只对原文页图成立，中译本本身是 40 页重排。
\item “affine level” $=k+h^\vee$；原文 §3.3 写 “level $k$” 是简写。
\end{itemize}

\section{公式与页码索引}
\label{app:equations}

本表只列正文实际引用或讨论过的公式，页码为 CMP 印刷页（351–399）。译本用页边注「原文 p.N」对应，PDF 页码为 40 页重排，不能直接换算。

\begin{center}
\small
\begin{longtable}{@{}p{2.1cm}p{6.0cm}p{1.4cm}p{4.6cm}@{}}
\toprule
编号 & 内容 & 页 & 在本书哪章 \\
\midrule
\endhead
(1.1) & 规范变换与曲率 $F_{ij}=[D_i,D_j]$ & 353 & 第 3 章 \\
(1.2) & Yang--Mills 作用量（依赖度规，被弃用） & 353 & 第 3 章 \\
(1.3) & Chern--Simons 作用量 & 354 & 第 3 章 \\
--- & 大规范变换 $L\to L+\mathrm{const}\cdot m$ & 354 & 第 3 章 \\
--- & 绕数归一化（$g^{-1}\dd g$ 的三次式） & 354 & 第 3 章 \\
(1.5) & Wilson 线 & 355 & 第 3、4 章 \\
(1.6) & 链环配分函数 & 355 & 第 3 章 \\
(2.16)(2.17) & 单圈修正 $k\to k+h^\vee$ & 359 & 第 5 章 \\
(2.25) & 大 $k$ 时 $c\to\dim G$ & 361 & 第 7 章 \\
(3.1) & $A_0=0$ 下的作用量（对时间一阶） & 367 & 第 4 章 \\
(3.2) & 泊松括号（系数 $4\pi/k$） & 367 & 第 4 章 \\
(3.3) & 高斯定律 $\epsilon^{ij}F^a_{ij}=0$ & 367 & 第 4 章 \\
(3.4) & 带标记点的高斯定律（算符值右端） & 371 & 第 4 章 \\
(3.5) & 大 $k$ 时 $\cH=\mathrm{Inv}(\otimes_i R_i)$ & 373 & 第 5 章 \\
(3.6) & $R\otimes R=\oplus_i E_i$（四点维数） & 374 & 第 5 章 \\
(4.20) & 标架反常相位 & 382 & 第 5 章 \\
(4.21) & $q=\exp\frac{2\pi\ii}{N+k}$ & 382 & 第 5、6 章 \\
(4.27)(4.28) & 手术律 & 385--386 & 第 6 章 \\
(4.36) & $N_{abc}$ 的定义（$S^2\times S^1$ 配分） & 388 & 第 6 章 \\
(4.41) & 标架反常 $Z\to Z\,e^{2\pi\ii(n-m)c/24}$ & 390 & 第 6 章 \\
(4.49) & 秩一性 & 393 & 第 6 章 \\
(4.51) & Verlinde 关键式（$S$ 对角化融合规则） & 393 & 第 6 章 \\
\bottomrule
\end{longtable}
\end{center}

三处已知译本笔误（引用前先看）：
\begin{enumerate}[leftmargin=2em]
\item 译本 §3.2 与 p.367--368 有三处把约束 (3.3) 误引为 (3.2)；读时一律按 (3.3)。
\item 原刊 (1.2) 的 Tr 归一化由脚注 1 给出；不要按 $-\frac14\int\sqrt g\,\Tr F^2$ 去补 $1/4$。
\item 早期资料曾说 p.353 漏 $1/4$，S1 证据官已判否：原刊照印，无 $1/4$、无负号。
\end{enumerate}

其它已核事实：
\begin{itemize}[leftmargin=1.6em]
\item 译本 40 页（重排），原文 49 页；公式编号沿用原文。
\item 原刊引言没有「本文解决第二个问题」这句话——那是译本导语框所加。
\item 本文是 1988 年 7 月 Swansea IAMP 报告的扩充版（p.351 脚注 *）；收稿 1988-09-12，修改 1988-09-27。
\end{itemize}


\end{document}
